% Part: counterfactuals
% Chapter: introduction
% Section: counterfactuals

\documentclass[../../../include/open-logic-section]{subfiles}

\begin{document}

\olfileid{cnt}{int}{cnt}

\olsection{પ્રતિતથ્યાત્મક શરતી વિધાનો}

અંગ્રેજીમાં ``જો \dots તો \dots'' પ્રકારની વાક્યરચનાનું એક અત્યંત
સામાન્ય અને મહત્વનું સ્વરૂપ \emph{to be} ક્રિયાપદના ભૂતકાલીન
સંભાવનાર્થી રૂપથી બને છે: ``જો \dots એવું થયું હોત, તો \dots એવું
થયું હોત.'' સામાન્ય રીતે આવા શરતી વિધાનનો પૂર્વપક્ષ અસત્ય, એટલે
હકીકતથી વિપરીત, હોય છે; તેથી તેમને \emph{પ્રતિતથ્યાત્મક શરતી
  વિધાનો} કહે છે. \emph{to be}ના સંભાવનાર્થી રૂપને કારણે તેમને
\emph{સંભાવનાર્થી શરતી વિધાનો} પણ કહે છે. તે \emph{નિવેદક} શરતી
વિધાનોથી જુદાં છે, જેમનું સ્વરૂપ ``જો \dots એવું છે, તો \dots એવું
છે'' હોય છે. પ્રતિતથ્યાત્મક અને નિવેદક શરતી વિધાનોની સત્ય થવાની
શરતો જુદી છે. એડમ્સનું જાણીતું ઉદાહરણ જુઓ:
\begin{quote}
  જો ઓસ્વાલ્ડે કેનેડીની હત્યા કરી ન હતી, તો બીજા કોઈએ કરી હતી.
  
  જો ઓસ્વાલ્ડે કેનેડીની હત્યા કરી ન હોત, તો બીજા કોઈએ કરી હોત.
\end{quote}
પહેલું વિધાન નિવેદક છે અને બીજું પ્રતિતથ્યાત્મક. પહેલું સ્પષ્ટપણે
સત્ય છે: આપણે જાણીએ છીએ કે પ્રમુખ જૉન એફ. કેનેડીની હત્યા
\emph{કોઈએ} કરી હતી; અને જો એ વ્યક્તિ લી હાર્વી ઓસ્વાલ્ડ ન હતી
(વૉરેન અહેવાલથી વિપરીત), તો કેનેડીની હત્યા બીજા કોઈએ કરી હતી. બીજું
વિધાન કંઈક જુદું કહે છે. તે દાવો કરે છે કે જો ઓસ્વાલ્ડે કેનેડીની હત્યા
ન કરી હોત---એટલે ડૅલસમાં ગોળીબાર ટળી ગયો હોત અથવા નિષ્ફળ ગયો
હોત---તો ત્યાર પછી ઇતિહાસ એવી રીતે આગળ વધ્યો હોત કે બીજી કોઈ
હત્યાનો પ્રયાસ સફળ થયો હોત. તે સત્ય હોય, તો કેનેડીનું મૃત્યુ
સુનિશ્ચિત કરવા શક્તિશાળી પરિબળોએ કાવતરું રચ્યું હોવું જોઈએ (જેમ
કેનેડીની હત્યા અંગે કાવતરું થયું હતું એવો સિદ્ધાંત માનનારાઓમાંના
અનેક લોકો માને છે).

\emph{નિવેદક} શરતી વિધાનને ભૌતિક શરતી વિધાન યોગ્ય રીતે રજૂ કરે છે
કે નહિ, તે હજુ ચર્ચાનો વિષય છે. ખાસ કરીને, અંગ્રેજીનાં નિવેદક શરતી
વિધાનોની સત્ય થવાની શરતો ભૌતિક શરતી વિધાન આપે છે એ વિચાર સાથે
સુસંગત રીતે તેના વિરોધાભાસોને ``સમજાવી'' શકાય કે નહિ, તે પ્રશ્ન છે.
તેનાથી વિપરીત, પ્રતિતથ્યાત્મક શરતી વિધાનોનું ભૌતિક શરતી વિધાન વડે
યોગ્ય પ્રતીકીકરણ થઈ શકતું નથી, એ નિર્વિવાદ છે. આ સ્પષ્ટ છે, કારણ કે
પ્રતિતથ્યાત્મક વિધાનના પૂર્વપક્ષો સામાન્ય રીતે અસત્ય હોવા છતાં,
અસત્ય પૂર્વપક્ષ ધરાવતાં બધાં પ્રતિતથ્યાત્મક વિધાનો સત્ય હોતાં નથી.
દાખલા તરીકે, તમે વૉરેન અહેવાલ માનતા હો અને કેનેડીની હત્યા માટે કોઈ
કાવતરું ન રચાયું હોય, તો એડમ્સનું પ્રતિતથ્યાત્મક
વિધાન આવું ઉદાહરણ છે.

કારણસંબંધી તર્કવિચારમાં પ્રતિતથ્યાત્મક શરતી વિધાનો મહત્વની ભૂમિકા
ભજવે છે: કારણસંબંધ વ્યક્ત કરવો એ તેમનો મુખ્ય ઉપયોગ છે. દાખલા તરીકે,
દીવાસળી ઘસવાથી તે સળગે છે; તેને ``જો આ દીવાસળી ઘસાત, તો તે સળગી
ઊઠત'' એમ કહી શકાય. ભૌતિક શરતી વિધાનો, અને સામાન્ય રીતે નિવેદક
શરતી વિધાનો, આ સંબંધ વ્યક્ત કરી શકતા નથી: ``દીવાસળી ઘસવામાં આવે છે
$\lif$ દીવાસળી સળગે છે'' એ વિધાન દીવાસળી ક્યારેય ઘસવામાં ન આવે
તો સત્ય છે, ભલે ઘસવામાં આવે ત્યારે શું થાય તે ગમે તે હોય. એથી પણ
વધુ વિચિત્ર વાત એ છે કે ``દીવાસળી ઘસવામાં આવે છે $\lif$ દીવાસળી
ફૂલોનો ગુચ્છો બની જાય છે'' એ વિધાન પણ દીવાસળી ક્યારેય ઘસવામાં ન
આવે તો સત્ય છે; પરંતુ ઘસવામાં આવે તો તે ચોક્કસ ફૂલોનો ગુચ્છો નહિ
બને.

પ્રતિતથ્યાત્મક શરતી વિધાનોનું યોગ્ય તર્કશાસ્ત્ર ચોક્કસ શું છે, તેની
હજુ ચર્ચા થાય છે. સ્ટાલનેકર અને લૂઈસે તેમનું પ્રભાવશાળી વિશ્લેષણ
આપ્યું. તેમના મતે ``જો~$S$ એવું થયું હોત, તો~$T$ એવું થયું હોત''
પ્રતિતથ્યાત્મક વિધાન ત્યારે અને માત્ર ત્યારે સત્ય છે જ્યારે વાસ્તવિક
જગતની સ્થિતિને સૌથી વધુ મળતી આવતી અને જ્યાં~$S$ સત્ય છે એવી
પ્રતિતથ્યાત્મક સ્થિતિમાં (``શક્ય જગતમાં'') $T$ સત્ય હોય. આને
``સત્તાલક્ષી'' વિશ્લેષણ કહે છે, કારણ કે તે શક્ય જગતોની
અસ્તિત્વમીમાંસાનો આધાર લે છે. બીજાં વિશ્લેષણો શરતી સંભાવનાઓ અથવા
માન્યતાઓના પુનરીક્ષણના સિદ્ધાંતો વાપરે છે. પ્રતિતથ્યાત્મક વિધાનોના
જુદાં જુદાં સૂચિત તર્કશાસ્ત્રોની સંખ્યા ઘણી છે. લૂઈસ--સ્ટાલનેકરનું
પણ એકમાત્ર તર્કશાસ્ત્ર નથી: સ્ટાલનેકર અને લૂઈસે સૌથી નજીકનાં શક્ય
જગતોના આધારે મળતાં આવતાં વર્ણનો આપ્યાં હોવા છતાં, તેમની અલગ
પૂર્વધારણાઓમાંથી જુદાં માન્ય અનુમાનો ફલિત થાય છે.

\end{document}
